\chapter{Electrostatics}\label{ch:electrostatics}

\section{Introduction}
Electrostatics occupies a distinctive position in geophysics. Like gravity and magnetics, it is governed by an elliptic potential-field equation in the steady state and can be formulated entirely in terms of scalar potentials and boundary-value problems. Unlike gravity, magnetics, and geothermics, however, most electrostatic measurements in geophysics rely on active sources, in which currents are deliberately injected into the Earth and the resulting electric potentials are observed. This distinction between passive and active methods is as important conceptually as the distinction between elliptic and parabolic governing equations.

In its steady-state form, electrostatics provides a natural bridge between potential-field methods and time-dependent electromagnetic techniques. Electrical potential, current flow, and conductivity can be treated using the same mathematical framework as other elliptic problems, while making explicit how material properties control field behavior. When time dependence is introduced, electrostatics transitions directly into electromagnetic diffusion, forming the conceptual foundation for transient EM and frequency-domain methods such as magnetotellurics.

This chapter introduces electrostatics with a deliberately limited scope. The goal is not to develop a full exploration methodology, but to establish the physical and mathematical ideas needed to understand how electrical methods connect to gravity, magnetics, geothermics, and ultimately to electromagnetic geophysics. In this sense, electrostatics functions as a conceptual bridge rather than a stand-alone field method.

This role should be understood in the context of the comparison of potential fields introduced at the end of the geothermics chapter, where electrostatics marks the transition from passive field observations to controlled, active-source probing of Earth structure.


\ntextbox{Key Concepts}{
    \begin{itemize}
        \item \textbf{Electrostatics is governed by an elliptic potential-field equation in the steady state.}
        Electrical potentials satisfy Poisson's or Laplace's equation, placing electrostatics on the same mathematical footing as gravity and magnetics under equilibrium conditions.

        \item \textbf{Most electrostatic measurements in geophysics are active-source techniques.}
        Currents are deliberately injected into the Earth, and observed potentials reflect the Earth's response to controlled forcing rather than natural fields.

        \item \textbf{Electrical conductivity is the primary material property controlling electrostatic responses.}
        Variations in lithology, porosity, fluid content, and temperature govern current flow and potential distributions.

        \item \textbf{Electrostatic fields respond instantaneously to changes in current distribution.}
        In the static limit there is no intrinsic temporal memory in the field, and observed potentials represent a snapshot of present-day conditions.

        \item \textbf{Electrostatics provides the steady-state limit of electromagnetic problems.}
        When time dependence is introduced, electrostatic formulations transition naturally to electromagnetic diffusion.

        \item \textbf{Self-potential is the primary passive electrostatic method in geophysics.}
        Natural electrical sources such as fluid flow, redox reactions, and thermal gradients generate measurable electric potentials without imposed currents.

        \item \textbf{Electrostatics clarifies the distinction between potential-field methods and EM geophysics.}
        It highlights which aspects of EM behavior arise from time dependence rather than from the underlying physics of electrical conduction.
    \end{itemize}
}{core}

\section{Mathematical symbols and units}
Electric potentials and fields
Definitions of electric potential, electric field, and current density

\section{Constitutive relations and material properties}
Ohm's law, conductivity, and geological controls on electrical properties

\section{Governing equations}
Poisson's and Laplace's equations for electrostatics; boundary-value formulation

\section{Active-source electrostatics}
Conceptual treatment of current injection, potential measurements, and common geometries

\section{Passive electrostatics: self-potential}
Natural sources of electric potentials and geological interpretations

\section{Connection to electromagnetic diffusion}
From steady-state electrostatics to time-dependent EM problems

\section{Relationship to other geophysical methods}
Comparison with gravity, magnetics, and geothermics; passive vs active methods

\section{Looking Ahead}
Transition to transient EM, frequency-domain methods, and magnetotellurics
This chapter is intentionally compact and conceptual, serving as a foundation for electromagnetic methods rather than a comprehensive treatment of electrical exploration.